Let $\mathcal{F}$ be a finite family of finite simple graphs.
Define $\mathrm{ex}(n;\mathcal{F})$ to be the maximum number of edges in an $n$-vertex graph containing no subgraph isomorphic to any $F\in\mathcal{F}$.

Assume $\mathcal{F}$ contains at least one bipartite graph.
The question is whether there always exists a \emph{bipartite} graph $G\in\mathcal{F}$ and a constant $C_{\mathcal{F}}>0$ such that for all $n$,
\[
\mathrm{ex}(n;G) \le C_{\mathcal{F}}\, \mathrm{ex}(n;\mathcal{F}).
\]
(Using Vinogradov notation: $\mathrm{ex}(n;G)\ll_{\mathcal{F}} \mathrm{ex}(n;\mathcal{F})$.)
